% 0_resumo.tex
\chapter*{Resumo}

O procurement assenta em documentos, e cada documento defeituoso que passa custa muito mais a corrigir a jusante do que a apanhar à entrada. Esta dissertação trata uma dessas barreiras na Hilti AG: a validação de cada Aviso de Expedição (ASN, do inglês \emph{Advanced Shipping Notice}) contra a respetiva Ordem de Compra (PO, do inglês \emph{Purchase Order}). Hoje essa verificação é um filtro binário de aceitação ou rejeição sobre campos estáticos; não faz validação semântica nem entre documentos, e quando falha arrasta o comprador para um ciclo de reconciliação manual de 30 a 90 minutos por defeito.

Os grandes modelos de linguagem são tentadores aqui, porque leem campos específicos de cada fornecedor e julgam termos contratuais com facilidade, mas são pouco fiáveis na aritmética que a tarefa também exige. Trabalhos recentes combinam-nos com verificação determinista, mas nenhum estudo comparou as principais arquiteturas híbridas frente a frente, deixando por medir os compromissos de exatidão, custo e latência que decidem qual construir. Esta dissertação fecha essa lacuna. Formaliza o padrão de \emph{enforcer}, uma camada determinista que detém a autoridade sobre o veredicto de todas as regras que um cálculo pode decidir, e apresenta o primeiro \emph{benchmark} das principais arquiteturas híbridas de validação numa única tarefa de 22 regras.

Enquadrada na metodologia de \emph{Design Science Research}, a solução é um pipeline de 7 etapas que substitui a barreira binária por uma decisão faseada de PASS, REVIEW ou REJECT. É avaliada sob 5 configurações que diferem na presença e na posição do \emph{enforcer}, sobre 67 casos sintéticos gerados de forma determinista e 32 casos reais cuja referência resulta do consenso de 3 analistas de procurement. Um único modelo de linguagem é mantido fixo em todas as configurações, sendo usado um segundo modelo de uma família diferente para testar a robustez, e as configurações são comparadas em concordância de disposição, exatidão por regra, custo e latência, com intervalos de confiança por \emph{bootstrap}.

O \emph{enforcer} separa as configurações. As 3 que o incluem atingem 0,88 de exatidão de disposição nos casos reais, contra 0,81 sem ele, uma diferença direcional nesta amostra pequena; nos 67 casos sintéticos, onde o contraste é estatisticamente fiável, essa diferença alarga para 20 pontos. A posição do \emph{enforcer} é uma escolha operacional e não de exatidão: decidir as regras deterministas antes da chamada ao modelo resulta na configuração ótima, igualando a melhor exatidão e recuperando mais regras, liderando no custo, com 0,073 EUR por ASN, e na latência, com 7,89 segundos de mediana. Um histórico do fornecedor nada acrescenta uma vez presente o \emph{enforcer}, e sob o segundo modelo a vantagem do \emph{enforcer} não só se mantém como aumenta, o indício mais claro de que o ganho vem da arquitetura e não do modelo. O trabalho contribui com um padrão arquitetural reutilizável, o primeiro \emph{benchmark} reportado das principais arquiteturas híbridas de validação numa única tarefa de 22 regras, e uma recomendação de implementação fundamentada em evidência para a validação de documentos de procurement.

\vspace{2em}
\noindent\textbf{Palavras-chave:} modelos de linguagem; pipeline híbrido LLM-determinista; padrão \emph{enforcer}; validação de documentos; procurement; \emph{design science research}.
